\section{Poisson distribution}

The goal of this section is to write a function for the Poisson probability distribution
\begin{equation}
P_\lambda(k) = \frac{\lambda^k e^{-\lambda}}{k!},
\end{equation}
that is accurate to 6 significant digits,
while only using 32-bit \texttt{numpy} datatypes.

Calculating the entire factorial $k!$ produces overflow for $k > 12$. 
This can be solved by changing the equation into a logarithm:
\begin{equation}
\begin{split}
\ln(P_\lambda(k)) & = k \cdot \ln(\lambda) - \lambda - \ln(k!) \\
& = k \cdot \ln(\lambda) - \lambda - \sum^{k}_{i=1} \ln(i)
\end{split}
\end{equation}
and returning $\exp(\ln(P_\lambda(k)))$.


To check the accuracy of our approximation, compare with the value given by \texttt{scipy.stats.poisson.ppf()}.
%The function \texttt{numpy.frexp()} decomposes a float into its fraction and its exponent. 
By looking at the difference between of our value and the fraction of \texttt{scipy}'s value, we can see from which digit they start to differ.
An example is given in Figure \ref{fig:question_1}, for $(\lambda,k)=(101,200)$.

%For $(\lambda,k) = (1,0), (5,10), (3,21)$ this works well enough.
%, with the fractional error staying well between the boundary of $\pm10^{-6}$.
%However, for $(\lambda,k) = (100,5), (101,200)$ the absolute fractional error is greater than $10^{-6}$, 
%and while $(\lambda,k) = (2.6,40)$ is within the tolerance, the fractional error increases for higher $k$.

\begin{figure}[h!]\label{fig:question_1}
  \centering
  \includegraphics[width=0.9\linewidth]{./plots/question1_101.png}
  \caption{
  Upper panel: Poisson distribution for $(\lambda,k) = (101,200)$. \\
  Lower panel: the fractional difference between the result from \texttt{Poisson\_distr()} and \texttt{scipy.stats.poisson.ppf()}.
  The horizontal grey lines indicate $\pm 10^{-5}$.
  If the reservoir trick isn't used (black line), the error jumps quite dramatically for $k>128$. When using the reservoir trick with a threshold of 32 (blue line), the error stays more constant.
  The vertical dashed line indicates $k$. 
  }
  \label{fig:fig1}
\end{figure}

If $k$ becomes greater than $128$, the error jumps quite dramatically.
This can be attributed to underflow.
For example, for $k=37$,
\begin{equation}
\begin{split}
\ln(37!) = \sum^{37}_{i=1} \ln(i) = \mathtt{99.33061}.
\end{split}
\end{equation}
However, if $k=38$, the extra $\mathtt{3.637586}$ gives us a total with 3 digits before the decimal point,
but the number of available significant digits stays the same, so we lose the information from the last digit:
\begin{equation}
\begin{split}
\ln(38!) = \ln(37!) + \ln(38) = \mathtt{102.9682},
\end{split}
\end{equation}
If we can make sure the running total never exceeds 100, we don't lose the information beyond the fifth decimal.
%37.0 99.33061 3.610918
%38.0 102.9682 3.637586

This is accomplished by putting the $k \cdot \ln(\lambda)$ term in a `reservoir',  and when the running total is about to exceed a set threshold, adding $k_\mathrm{step} \cdot \ln(\lambda)$ to it (and reducing the reservoir accordingly).

While the example above uses base 10, of course a computer uses base 2, 
so when choosing a suitable threshold, it makes sense to try powers of 2. 
Down to 64, each power of 2 performs better than the one above.
For thresholds below 16, the accumulation of machine error starts to worsen our results.
There is no apparent difference between the results of thresholds 16, 32 and 64, so choose 32 as the middle ground.





All this is implemented in the following script:
\lstinputlisting{question1.py}

Output:
\lstinputlisting{question1_output.txt}

